%=====================================================================
\section{Related Work}
\label{sec:related}
We organise the literature by question rather than paper by paper; the subsections
follow Parts I--IV of the Introduction. General frameworks and reviews are given by
\citet{holme2012temporal,holme2015modern} and the book of \citet{masuda2016guide}.
A central point of these reviews is that the basic unit is a \emph{contact}
(two nodes and a time), and that connectedness is no longer transitive: if $a$ meets
$b$ and later $b$ meets $c$, a path from $a$ to $c$ exists, but not the other way
round. Hence a temporal network cannot be turned into a simple graph without losing
information or changing what a node means, and many static-network methods need
non-trivial generalisations \citep{holme2015modern}.

\subsection{Representing time: aggregation, snapshots and windows}
\label{sec:rw-representation}
Temporal networks can be stored as contact sequences, interval graphs or sequences of
snapshots \citep{holme2012temporal}. Lossy representations such as weighted aggregates or reachability graphs reduce the
data to a static object; lossless ones keep all contacts \citep{holme2015modern}.
Each choice involves a trade-off: snapshots
make static tools applicable but require a window length, whereas the full event
sequence keeps all information. The effect of this choice has been quantified for
communication data: window size and placement change the structure of the aggregated
network \citep{krings2012effects}, and coarser temporal resolution distorts
properties of the time-varying network in a way that depends on the system
\citep{ribeiro2013quantifying}. This motivates \textbf{H1}: aggregation and window
choice change structural conclusions.

\subsection{Temporal activity patterns}
\label{sec:rw-activity}
Human activity is bursty: inter-event times are heavy-tailed rather than exponential
\citep{barabasi2005origin}, and consecutive intervals are correlated, which
\citet{goh2008burstiness} capture with a burstiness coefficient and a memory
coefficient. Social ties are also repeated and persistent over long periods
\citep{saramaki2014persistence}, while the set of active links turns over. Temporal motifs
\citep{kovanen2011temporal,paranjape2017motifs} and activity-driven models
\citep{perra2012activity} describe and generate such patterns. This gives
\textbf{H2} (activity is bursty and has memory) and \textbf{H3} (ties are repeated and
persistent while the network turns over).

\subsection{Structural evolution: growth, ageing and communities}
\label{sec:rw-evolution}
Growth is classically explained by preferential attachment
\citep{barabasi1999emergence}, whose strength can be measured directly from
growth data \citep{jeong2003measuring,newman2001clustering}. Real systems usually
show ageing, where old nodes attract fewer new links
\citep{dorogovtsev2000aging}, and in citations this is captured by a relevance
decay \citep{wang2013quantifying}. Growing networks also densify, with the number of
edges growing superlinearly in the number of nodes
\citep{leskovec2005graphs,leskovec2007graph}; node and link deletion alter these
laws \citep{moore2006exact}. This gives \textbf{H4}: growth follows preferential
attachment with ageing, and networks densify.

At the mesoscale, multilayer modularity couples partitions across time
\citep{mucha2010community}, communities can be tracked across snapshots through
matching \citep{greene2010tracking} and change can be mapped between partitions
\citep{rosvall2010mapping}; \citet{rossetti2018community} survey the field. Per-snapshot
partitions are typically obtained with Louvain \citep{blondel2008fast} or Leiden
\citep{traag2019louvain}. This gives \textbf{H5}: communities reorganise over time and
a static partition hides this.

\subsection{Time and network function}
\label{sec:rw-function}
Because interactions occur in sequence, only time-respecting paths can carry
information: reachability in contact sequences is much lower than in the aggregated
graph \citep{holme2005network}, and path lengths and centrality change accordingly
\citep{pan2011path,kim2012temporal}. Burstiness and temporal correlations further
slow down spreading even when the static topology is unchanged
\citep{karsai2011small}. These results give \textbf{H6} (time ordering restricts
reachability and changes centrality rankings) and \textbf{H7} (burstiness and
temporal correlations slow down spreading).

\subsection{Randomized reference models}
\label{sec:rw-null}
Randomization is more versatile for temporal than for static networks, where it is
essentially limited to degree-preserving link shuffling, because one can switch off
individual kinds of correlation \citep{holme2015modern}. Randomized reference models act as hypothesis tests: each one preserves some
properties of the data and destroys others, so differences between data and model
reveal which feature matters \citep{gauvin2022randomized}. We use them in Parts
II--IV (\cref{sec:null}).

\paragraph{Time scales and network types.}
Temporal networks are usually defined as systems whose evolution is not much slower
than the processes running on them, and a single system contains many time scales
(seconds, hours, days) \citep{holme2015modern}. Contact and communication data are
event sequences with such a hierarchy, whereas citation networks are acyclic by
construction and all out-links of a paper appear at one instant
\citep{holme2015modern}. This is why we expect, and test, different sensitivity to time
across network types.

\paragraph{Gap.}
Most studies examine one dataset or one method. Few compare, within a single
framework, how sensitive different \emph{types} of networks are to temporal
information. The hypotheses distilled above (\cref{tab:hypotheses}) allow such a
comparison.
